\section{Conclusion and Future Work}
\label{sec:conclusion}

We presented a hybrid GNN-LLM fusion architecture for Ethereum smart-contract
fraud detection that combines a two-layer GraphSAGE network operating on an
eight-dimensional node feature vector with a frozen CodeBERT encoder over
verified Solidity source code.
A reproducible four-phase pipeline---BigQuery transaction retrieval, DuckDB
integration, PyTorch Geometric graph construction, and CodeBERT tokenisation---
produces a dense, label-rich graph of \num{85680} nodes and \num{239742}
transactions from which \num{304} contracts carry both verified source and an
explicit fraud/benign label.

On the held-out test split, the concatenation-based fusion model attains
\textbf{93.6\% accuracy}, a \textbf{fraud F1-score of 0.909},
\textbf{AUC-ROC of 0.961}, and \textbf{AUC-PR of 0.959}---achieving the
highest F1-Fraud and MCC (0.861) across all seven evaluated models.
A comprehensive ablation study compared GraphSAGE-only (F1\,=\,0.636),
GCN-only (F1\,=\,0.667), GAT-only (degenerate on small graphs),
CodeBERT-only (F1\,=\,0.514), a Random Forest baseline (F1\,=\,0.889),
and an attention-based fusion variant (F1\,=\,0.625) under identical
conditions.  The Random Forest's strong AUC (0.988) validates the
discriminative power of our eight-dimensional node features, while the
fusion model's superior F1 and MCC confirm that code semantics provide
complementary signal unavailable to any purely graph-based approach.

The modular codebase released alongside the paper replaces an error-prone
monolithic notebook with independently testable components, and the two
explicit seed arrays make the dataset definition fully auditable and
extensible.

\textbf{Future work} should address the following priorities:

\begin{enumerate}
    \item \textbf{Evaluation hardening.}  Replace the random split with
    temporal and deployer-disjoint splits; report mean $\pm$ std over at
    least five seeds; evaluate on a shared public Ethereum phishing
    benchmark to enable direct comparison with TTAGN~\cite{li2022ttagn},
    BERT4ETH~\cite{hu2023bert4eth}, and TLMG4Eth~\cite{sun2025tlmg4eth}.

    \item \textbf{Dataset expansion.}  The fraud F1 ceiling is set by the
    small absolute count of labelled fraud contracts with verified source.
    Curating additional DeFiHackLabs victim contracts, incorporating
    bytecode decompilation for unverified contracts, and extending the
    seed list to post-2026 exploits would likely improve recall more than
    architectural tuning.

    \item \textbf{Parameter-efficient LLM fine-tuning.}  LoRA
    adapters~\cite{hu2022lora} would allow CodeBERT to specialise on
    Solidity vulnerability patterns without exceeding the memory budget
    of a consumer GPU.

    \item \textbf{Temporal graph modelling.}  Replacing the static
    GraphSAGE encoder with a temporal graph network that encodes edge
    timestamps and transaction sequences would let the model capture
    the bootstrapping phase of a fraud contract---the brief window
    before funds are drained---which is invisible to the current
    static-graph representation.

    \item \textbf{Cross-chain generalisation.}  Applying the pipeline
    to BNB Chain, Polygon, and Arbitrum would test whether the fusion
    architecture generalises across EVM-compatible networks with
    different transaction-fee structures and DeFi ecosystems.
\end{enumerate}

The combination of these directions, together with the reproducible
pipeline and curated seed arrays released with this paper, provides a
solid foundation for the next generation of multi-modal Ethereum fraud
detectors.
